\documentclass[11pt]{article}

\usepackage[margin=1in]{geometry}
\usepackage{amsmath,amssymb,amsthm}
\usepackage{booktabs}
\usepackage{graphicx}
\usepackage{hyperref}
\usepackage{xcolor}
\usepackage{microtype}

\newtheorem{theorem}{Theorem}
\newtheorem{proposition}{Proposition}
\newtheorem{corollary}{Corollary}

\newcommand{\R}{\mathbb{R}}
\newcommand{\E}{\mathbb{E}}
\newcommand{\norm}[1]{\lVert #1\rVert}
\newcommand{\mwg}{\textnormal{\textsc{MWG-EW}}}
\newcommand{\traffic}{\mathcal{T}}
\newcommand{\loss}{\mathcal{L}}

\title{\bf Ephemeral Meta-Weights for Bandwidth-Aware Transformer Feed-Forward Networks}
\author{Anonymous Authors}
\date{May 2026}

\begin{document}
\maketitle

\begin{abstract}
Transformer inference and training increasingly run into a weight-bandwidth
limit: feed-forward network (FFN) projections are large, reused at every token,
and repeatedly streamed from off-chip memory. We study \mwg{}, a neural
architecture that replaces selected persistent FFN matrices with conditional
low-rank meta-weights. A small generator produces rank-$r$ descriptors for the
current layer and token group; the target runtime consumes those descriptors
with an associative $(XU_r)V_r$ schedule and a no-write-back lifecycle. In the
transformed block, the dense FFN matrix is no longer stored or synchronized as
a trainable object, while the present prototype measures an unfused
associative path and still requires fused-kernel counter validation.

This paper makes three contributions. First, we define a generator-conditioned
Transformer FFN that preserves differentiability while exposing rank, token
reuse, and basis-bank size as architecture knobs. Second, we give an
infrastructure contract for ephemeral descriptor execution and derive compute,
bandwidth, memory, latency, and gradient-communication scaling laws. Third, we
provide a distributed evaluation harness and report wall-clock latency, traffic
estimates, communication volume, checkpoint-backed approximation quality, and
small end-to-end perplexity gates on a patched Qwen2.5-1.5B-Instruct model.
On an eight-accelerator ASI1 Ascend 910B2 prototype, the current implementation
at $r=128$ reduces estimated descriptor bytes by 24.9x and trainable FFN state
by 18.4x, improves the measured training step by $1.542\pm0.196$x, and reduces
measured all-reduce latency by 5.4x. The fastest measured point, $r=256$,
reaches $1.572\pm0.165$x training speedup; the lower-rank $r=64$ point gives
49.8x estimated descriptor-byte reduction. A checkpoint-backed Qwen2.5-1.5B FFN
distillation experiment shows that short low-rank distillation improves over
static SVD but remains far from dense-output fidelity on early layers. A
quality-gated ASI3 probe identifies layer 16 as the least brittle replacement
candidate, while a two-layer patch that looked strong on a 4080-token diagnostic
gate does not generalize to a more diverse text stress set. Direct LM-CE
calibration improves layer 16 on the original held-out instruction slice, but a
new manifest-driven broad validation still degrades GSM8K, MBPP, and Alpaca-tail
suites, and a deployable router stays close to dense only when it patches a very
small token fraction. The central claim is therefore a feasibility case and
negative validation study rather than a final deployment claim: learned
conditional generation, eligibility selection, and a no-write-back operand
lifecycle can make parameter movement an architectural design variable, provided
future training, fused kernels, and hardware counters close the current quality
and traffic-evidence gaps.
\end{abstract}

\section{Introduction}

The standard Transformer block stores FFN projections as persistent dense
matrices and streams them during every forward pass. For autoregressive
decoding, the arithmetic per token can be small relative to the bytes fetched
from high-bandwidth memory. For distributed training, the same persistent
matrices induce equally persistent gradient synchronization traffic. This paper
asks whether a Transformer can improve this operating point by changing the
lifecycle of FFN weights rather than only changing their datatype.

\mwg{} turns a subset of FFN weights into generated, short-lived descriptors.
The descriptor is not a compressed checkpoint artifact; it is an operand
created by a neural generator, consumed by a kernel, and then overwritten. The
architecture is therefore both algorithmic and infrastructural: the learning
rule determines which descriptors are useful, while the runtime determines
whether the descriptors actually avoid off-chip traffic.

This distinction matters because many efficient-model methods still leave the
same lifecycle in place: a parameter tensor is stored, loaded, differentiated,
and synchronized as a durable object. \mwg{} studies the complementary question:
what happens if a large projection is a transient computation product rather
than a resident model object? If this view holds under full training and quality
validation, it suggests a new axis for scaling Transformers: not only fewer
parameters or fewer bits per parameter, but shorter parameter lifetimes.

The proposal is close in spirit to hypernetworks \cite{ha2017hypernetworks},
low-rank adaptation \cite{hu2022lora}, and IO-aware kernels such as
FlashAttention \cite{dao2022flashattention}, and it is complementary to
attention-side KV-cache compression such as Multi-head Latent Attention
\cite{deepseek2024v2}. Hypernetworks usually materialize generated weights as
ordinary model tensors. LoRA stores low-rank factors as persistent adapter
parameters. FlashAttention optimizes attention activation traffic, while MLA
compresses attention keys and values during autoregressive decoding. \mwg{}
instead treats generated FFN factors as volatile kernel operands.

\paragraph{Contributions.}
We make the following contributions:
\begin{itemize}
  \item A Transformer FFN replacement in which a learned generator emits
  conditional rank-$r$ descriptors for the up, gate, and down projections.
  \item A no-write-back descriptor lifecycle that connects the neural
  architecture to memory traffic and distributed synchronization.
  \item Scaling laws for parameter bytes, descriptor bytes, latency in the
  memory-bound regime, and all-reduce volume.
  \item A distributed benchmark that records latency, memory-traffic estimates,
  and communication evidence in machine-readable form.
  \item A negative static-SVD baseline and end-to-end patch tests showing why
  the method should be judged as conditional generation, eligibility
  selection, patch-aware calibration, and lifecycle control rather than
  ordinary low-rank compression.
  \item A manifest-driven broad-validation and deployable-router study showing
  that the current layer-16 checkpoint is not yet broadly quality preserving.
\end{itemize}

\section{Architecture}

\subsection{Baseline Gated FFN}

A gated Transformer FFN maps hidden states $X\in\R^{B\times S\times d}$ to
\begin{equation}
Y = \left(\phi(XW_{\mathrm{gate}})\odot XW_{\mathrm{up}}\right)W_{\mathrm{down}},
\end{equation}
where $W_{\mathrm{up}}, W_{\mathrm{gate}}\in\R^{d\times m}$,
$W_{\mathrm{down}}\in\R^{m\times d}$, and $\phi$ is typically SiLU. With fp16
weights, one block stores approximately $6dm$ bytes when $m$ dominates $d$.

\subsection{Ephemeral Meta-Weight FFN}

\mwg{} replaces each projection matrix by a conditional descriptor pair:
\begin{equation}
(U^{(p)}_r,V^{(p)}_r) = G_\theta(c_t,\ell,p),\qquad
XW_p \approx (XU^{(p)}_r)V^{(p)}_r,
\end{equation}
where $p\in\{\mathrm{up},\mathrm{gate},\mathrm{down}\}$, $c_t$ is a token-group
context, $\ell$ is the layer id, $U_r\in\R^{d_{\mathrm{in}}\times r}$, and
$V_r\in\R^{r\times d_{\mathrm{out}}}$. The FFN becomes
\begin{equation}
\widetilde{Y} =
\left(\phi((XU^{(g)}_r)V^{(g)}_r)\odot (XU^{(u)}_r)V^{(u)}_r\right)
U^{(d)}_rV^{(d)}_r.
\end{equation}

We use a basis-bank generator:
\begin{align}
\alpha &= \mathrm{softmax}(A_\ell h_\theta(c_t)),\\
U_r^{(p)} &= \sum_{k=1}^{K}\alpha_k U_{k}^{(p)} + \Delta U_\theta^{(p)}(c_t,\ell),\\
V_r^{(p)} &= \sum_{k=1}^{K}\alpha_k V_{k}^{(p)} + \Delta V_\theta^{(p)}(c_t,\ell).
\end{align}
The bank gives stable reusable directions; the residual head gives
input-conditioned adaptation. Practical kernels may generate only rank-channel
scales or tile-local residuals when full descriptor synthesis is too expensive.

\subsection{Descriptor Lifecycle}

The runtime contract is simple and testable:
\begin{enumerate}
  \item Read token context $c_t$ and layer id $\ell$.
  \item Generate descriptor factors $(U_r,V_r)$ for each FFN projection.
  \item Store descriptor tiles in an on-chip scratch region when the target
  accelerator exposes such memory; otherwise keep the same volatile lifecycle
  at the closest available cache level.
  \item Compute each projection with the associative schedule $(XU_r)V_r$.
  \item Release or overwrite each scratch tile before the next token group.
  \item Synchronize gradients only for generator and basis-bank parameters.
\end{enumerate}

The important distinction is not only low rank. A low-rank factor stored as a
long-lived model tensor still creates persistent parameter traffic. The
ephemeral contract allows a runtime to avoid writing generated descriptors back
to off-chip memory and allows distributed training to synchronize only the
generator-side state.

\section{Analysis}

\subsection{Traffic Scaling}

For a gated FFN with three projections and fp16 elements, dense weight traffic
per block is
\begin{equation}
\traffic_{\mathrm{dense}} = 6dm.
\end{equation}
Ignoring the small generator context, rank-$r$ descriptor traffic is
\begin{equation}
\traffic_{\mathrm{mwg}} = 2r(d+m) + 2r(d+m) + 2r(m+d)
  = 6r(d+m).
\end{equation}
The idealized traffic ratio is therefore
\begin{equation}
\Gamma(r)=\frac{\traffic_{\mathrm{mwg}}}{\traffic_{\mathrm{dense}}}
= \frac{r(d+m)}{dm}.
\end{equation}
At $d=4096$, $m=14336$, and $r=128$, this gives
$\Gamma=0.0402$, or a 24.9x reduction in descriptor bytes.

\begin{proposition}[Memory-bound speedup condition]
Assume per-block latency can be approximated by
$L=L_0+\kappa\traffic$ with fixed launch/compute term $L_0$ and bandwidth
coefficient $\kappa>0$. Then \mwg{} improves latency whenever
$\traffic_{\mathrm{mwg}}<\traffic_{\mathrm{dense}}$, and the asymptotic speedup
approaches $1/\Gamma(r)$ as $L_0$ becomes small relative to memory traffic.
\end{proposition}

\subsection{Gradient Communication}

In data-parallel training, dense FFN gradients require all-reduce traffic
proportional to $6dm$ bytes per block. \mwg{} synchronizes generator and
descriptor-bank parameters, whose leading term is $6r(d+m)$ plus the small MLP
state. The same ratio $\Gamma(r)$ controls the dominant gradient-byte reduction.
This is orthogonal to optimizer sharding and tensor parallelism: those systems
can still be used on the remaining generator parameters.

\subsection{Computational Complexity}

Let $T=BS$ be the number of tokens processed by one FFN call and let $q$ be the
token-group reuse factor for a generated descriptor. The dense gated FFN
performs three large matrix products and has leading compute
\begin{equation}
\mathcal{C}_{\mathrm{dense}} = \Theta(3Tdm)
\end{equation}
multiply-accumulate operations. The associative \mwg{} path computes
$XU_r$ and then $ZV_r$ for the up, gate, and down projections:
\begin{equation}
\mathcal{C}_{\mathrm{mwg,assoc}} = \Theta(3Tr(d+m)).
\end{equation}
If a basis bank with $K$ entries is explicitly blended once per token group,
the descriptor-generation term is
\begin{equation}
\mathcal{C}_{\mathrm{gen}} =
\Theta\!\left(3\left\lceil T/q\right\rceil K r(d+m) +
\left\lceil T/q\right\rceil C_G\right),
\end{equation}
where $C_G$ is the small context-network cost. Kernels that generate only
rank-channel scales or tile-local residuals replace the $Kr(d+m)$ term by the
corresponding tile size. Thus the leading compute ratio, ignoring amortized
generation, is again $\Gamma(r)=r(d+m)/(dm)$. The method is attractive when
$r \ll \min(d,m)$ and $q$ is large enough to amortize generation; it can be
unattractive when descriptors are regenerated for very small token groups or
when unfused factor matmuls lose accelerator utilization.

\subsection{Working-Set Complexity}

The persistent FFN parameter state is $\Theta(3dm)$ elements per layer. A
fully materialized rank-$r$ descriptor set is only $\Theta(3r(d+m))$ elements,
and a tiled kernel can reduce the live scratch working set further to
$\Theta(r(d_{\mathrm{tile}}+m_{\mathrm{tile}}))$. These are storage-footprint
and traffic estimates. They are not claims that every descriptor or activation
is resident in on-chip SRAM at once; the no-write-back contract says that
generated descriptor tiles should be consumed at the closest available
scratch/cache level and not committed as durable off-chip model weights.

\subsection{Approximation}

Let $\widetilde{W}_p=U^{(p)}_rV^{(p)}_r$ be the generated approximation of
$W_p$. If $\norm{W_p-\widetilde{W}_p}_F\le \epsilon_p$, then for bounded
activations and $L_\phi$-Lipschitz nonlinearity,
\begin{equation}
\norm{Y-\widetilde{Y}}_F
\le C_1 L_\phi\norm{X}_F(\epsilon_{\mathrm{gate}}+\epsilon_{\mathrm{up}})
  + C_2\norm{X}_F\epsilon_{\mathrm{down}},
\end{equation}
where $C_1$ and $C_2$ depend on activation bounds inside the FFN. This connects
rank selection to output distortion, while leaving room for end-to-end
distillation to improve beyond static SVD.

\subsection{Static Low-Rank Is Not the Claim}

To separate \mwg{} from ordinary compression, we analyze a static truncated-SVD
baseline on 18 FFN matrices from a Qwen-1.5B-family checkpoint. Table
\ref{tab:svd} reports mean retained singular-value energy and reconstruction
cosine. At ranks that are attractive for bandwidth, static SVD preserves too
little FFN structure to be a complete solution. This is evidence for the
generator side of \mwg{}: the low-rank factors should be conditional operands,
not merely fixed compressed weights.

\begin{table}[t]
\centering
\caption{Static SVD baseline on 18 pretrained FFN matrices. Low ranks are
excellent for bytes but weak as fixed approximations, motivating learned
conditional generation.}
\label{tab:svd}
\begin{tabular}{rrrr}
\toprule
Rank $r$ & Mean energy & Min energy & Mean cosine \\
\midrule
32 & 7.17\% & 4.27\% & 0.266 \\
64 & 11.98\% & 8.16\% & 0.345 \\
128 & 20.12\% & 15.46\% & 0.448 \\
256 & 33.79\% & 28.65\% & 0.582 \\
512 & 55.61\% & 50.75\% & 0.748 \\
1024 & 85.13\% & 82.42\% & 0.925 \\
\bottomrule
\end{tabular}
\end{table}

\section{Systems Realization}

An implementation needs three pieces:
\begin{itemize}
  \item \textbf{Generator placement.} The generator should run adjacent to the
  FFN kernel and emit descriptor tiles directly into scratch memory.
  \item \textbf{Associative execution.} The kernel computes $XU_r$ followed by
  $ZV_r$, fusing activation and gate operations where possible.
  \item \textbf{Lifecycle enforcement.} Descriptor buffers are marked volatile;
  runtime instrumentation checks that descriptor write-back does not occur.
\end{itemize}

The benchmark in this folder measures the associative path and the distributed
communication effect. A production kernel should additionally expose hardware
counters for off-chip descriptor reads and writes.

\section{Experimental Protocol}

\paragraph{Execution setup.}
The large-scale reliability run uses the ASI1 Huanxin environment with eight
visible Ascend 910B2 NPUs. We evaluate one, two, four, and eight data-parallel
ranks by setting \texttt{ASCEND\_RT\_VISIBLE\_DEVICES} to the corresponding
device set and launching PyTorch distributed execution with HCCL. Each
configuration is repeated three times. The raw artifact records platform
metadata for reproducibility, but the method and results below are reported in
terms of model geometry, rank, latency, state size, estimated traffic, and
collective communication.

\paragraph{Geometry.}
The large settings use an 8B-scale FFN block with $d=4096$, $m=14336$, decode
shape $B=1,S=128$, and training shape $B=2,S=512$. The reliability/scaling
sweep reported below uses representative ranks $r\in\{64,128,256\}$.

\paragraph{Validation path.}
The benchmark uses a small smoke geometry $(d,m)=(512,1408)$ with short
iteration counts before running the 8B-scale geometry. This keeps the expensive
distributed sweep separate from basic checks of output writing, tensor shapes,
and collective communication.

\paragraph{Metrics.}
We report forward latency, backward latency, throughput, trainable parameter
bytes, estimated descriptor traffic, dense-to-\mwg{} speedup, all-reduce
latency, and all-reduce bytes.

\paragraph{Checkpoint quality and end-to-end probes.}
To avoid treating systems speed as a proxy for model quality, we separately
distill the FFN output of layer 0 from a local Qwen2.5-1.5B-Instruct checkpoint
stored in the ASI1 file system. The teacher uses the checkpoint
\texttt{gate\_proj}, \texttt{up\_proj}, and \texttt{down\_proj} matrices with
$d=1536$ and $m=8960$. We compare static truncated SVD, a trainable persistent
low-rank FFN, and the same low-rank factors with the \mwg{} rank-scale
generator. The reported run uses eight NPUs, fp32 training for numerical
stability, ranks $64$ and $128$, 16 short distillation steps, and four
evaluation batches. We then run a stronger ASI3 quality path on the same
checkpoint: text-activation distillation, expert-residual \mwg{} students,
saved student checkpoints, and direct replacement of selected Qwen MLP modules
followed by causal-language-model perplexity evaluation. This second probe is
still small, but it is end-to-end: the patched model is evaluated through the
full Transformer, not only by block MSE.

\paragraph{Validation controls.}
After auditing an invalid recovery artifact that evaluated on the tiny default
calibration text stream, all held-out-named recovery and broad-validation
launchers now require explicit text inputs. The ASI3 recovery path refuses
\texttt{DEFAULT\_TEXTS} fallback through a \texttt{--require-texts} guard, and a
manifest-driven broad-validation harness rejects empty or undersized corpora
before launching PPL evaluation. These controls are procedural rather than new
quality evidence, but they are important for reproducibility: held-out results
below are only comparable when the text source and token count are explicit.

\paragraph{Deployable selective routing protocol.}
The next algorithmic validation step is a non-oracle dense-fallback policy for
the best current single-layer patch. We added a router evaluation harness that
uses one explicit text split to label examples by dense-vs-patched loss delta,
fits a small ridge regressor from features available before the target FFN
executes, and evaluates thresholded patch fractions on a disjoint held-out
split. The features are summary statistics of the layer-16 pre-FFN hidden state
and sequence occupancy, not evaluation-set dense-vs-patched losses or
post-hoc perplexity comparisons. Its report includes always-dense,
always-patched, deployable routed, and oracle-frontier rows. This protocol is
intended to test whether \mwg{} can become a selective runtime optimization
rather than an indiscriminate layer replacement; it does not add a positive
claim until broad held-out corpus files and multiple seeds have been run.

\paragraph{Artifact trail.}
The reported ASI1 systems run is stored under the historical artifact name
\texttt{ai3\_reliability\_scaling\_20260518T072507Z}. It contains 12 raw JSON
benchmark artifacts: three repeats each for world sizes 1, 2, 4, and 8, plus an
aggregate JSON/Markdown summary. The local repository keeps the pulled summary,
status file, and log under \texttt{results/pulled\_asi1/}; the same experiment
code and paper artifacts are mirrored to the project S3 folder.

\section{Numerical Evidence}

\paragraph{Evidence status.}
The results below were measured on 2026-05-18 on ASI1 using all eight visible
NPUs for the largest setting, \texttt{torch.distributed.run}, HCCL, and fp16
tensors. The benchmark uses the 8B-scale FFN geometry
$(d,m)=(4096,14336)$, decode shape $(1,128,4096)$, and training shape
$(2,512,4096)$. Each row reports the mean and standard deviation over three
independent launches. Each launch used three warmup iterations, ten forward
timing iterations, three training iterations, and five all-reduce iterations.
These measurements establish systems behavior for the associative low-rank
path and distributed communication. They do not yet establish final end-to-end
language-model quality.

\begin{table}[t]
\centering
\caption{Reliability and scaling of 8B-scale FFN training latency. Each cell is
mean$\pm$std over three launches on ASI1. The current prototype uses unfused
associative execution, not yet a fused no-write-back kernel.}
\label{tab:scaling}
\begin{tabular}{rrrrr}
\toprule
World & Rank & Train ms & Train speedup & Train tok/s \\
\midrule
1 & -- & $4.9742\pm0.0239$ & 1.000x & 0.206M \\
1 & 64 & $3.3785\pm0.1928$ & $1.475\pm0.075$x & 0.304M \\
1 & 128 & $3.2545\pm0.1694$ & $1.531\pm0.071$x & 0.315M \\
1 & 256 & $3.2879\pm0.1327$ & $1.514\pm0.055$x & 0.312M \\
2 & -- & $5.8624\pm0.7332$ & 1.000x & 0.353M \\
2 & 64 & $3.9907\pm0.1104$ & $1.472\pm0.211$x & 0.513M \\
2 & 128 & $4.7966\pm1.5295$ & $1.268\pm0.214$x & 0.453M \\
2 & 256 & $3.8896\pm0.0662$ & $1.507\pm0.181$x & 0.527M \\
4 & -- & $6.1351\pm1.1496$ & 1.000x & 0.682M \\
4 & 64 & $4.1078\pm0.1168$ & $1.497\pm0.303$x & 0.998M \\
4 & 128 & $4.4547\pm0.6583$ & $1.412\pm0.423$x & 0.932M \\
4 & 256 & $4.1109\pm0.0797$ & $1.496\pm0.312$x & 0.997M \\
8 & -- & $6.4874\pm0.8606$ & 1.000x & 1.279M \\
8 & 64 & $4.8044\pm1.1034$ & $1.373\pm0.182$x & 1.760M \\
8 & 128 & $4.2047\pm0.0299$ & $1.542\pm0.196$x & 1.948M \\
8 & 256 & $4.1222\pm0.2432$ & $1.572\pm0.165$x & 1.992M \\
\bottomrule
\end{tabular}
\end{table}

\begin{table}[t]
\centering
\caption{State size, descriptor traffic, and measured all-reduce latency for
the same ASI1 reliability run. AllReduce latency cells are mean ms over three
launches for the corresponding world size.}
\label{tab:statecomm}
\begin{tabular}{rrrrrrr}
\toprule
Rank & Params MiB & Descriptor MiB & Traffic red. & AR 2 & AR 4 & AR 8 \\
\midrule
-- & 336.000 & 336.000 & 1.000x & 19.6627 & 11.4896 & 6.9101 \\
64 & 11.127 & 6.750 & 49.778x & 1.2599 & 1.1655 & 1.1549 \\
128 & 18.254 & 13.500 & 24.889x & 1.6772 & 1.4632 & 1.2786 \\
256 & 32.507 & 27.000 & 12.444x & 2.5398 & 1.8239 & 1.4182 \\
\bottomrule
\end{tabular}
\end{table}

\begin{table}[t]
\centering
\caption{Checkpoint-backed FFN approximation on ASI1 using Qwen2.5-1.5B layer
0. Lower relative MSE and higher cosine are better. The short eight-NPU
distillation run improves over static SVD but shows that the current
rank-scale generator does not yet recover dense FFN fidelity at bandwidth-
attractive ranks.}
\label{tab:quality}
\begin{tabular}{lrrrrr}
\toprule
Method & Rank & Rel. MSE & Cosine & Params MiB & Traffic red. \\
\midrule
Static SVD & 64 & 0.989094 & 0.133395 & -- & 20.488x \\
Persistent low-rank & 64 & 0.966727 & 0.225573 & 9.096 & 20.488x \\
\mwg{} rank-scale & 64 & 0.966702 & 0.225632 & 9.096 & 20.488x \\
Static SVD & 128 & 0.971749 & 0.187598 & -- & 10.244x \\
Persistent low-rank & 128 & 0.942873 & 0.257047 & 17.066 & 10.244x \\
\mwg{} rank-scale & 128 & 0.942804 & 0.257052 & 17.066 & 10.244x \\
\bottomrule
\end{tabular}
\end{table}

\begin{table}[t]
\centering
\caption{Diagnostic end-to-end Qwen2.5-1.5B perplexity gates for patched
\mwg{} FFN layers on ASI3. Lower PPL ratio is better. The 1020--4080-token rows
are not validation claims; they show why early-layer and multi-layer policies
must be rejected or re-tested on held-out text.}
\label{tab:ppldiag}
\begin{tabular}{p{0.40\linewidth}rrrr}
\toprule
Patch policy & Rank & Tokens & Dense PPL & Patched PPL / ratio \\
\midrule
Layers 0,4,8,12, independent & 128 & 1020 & 314.44 & 2.18M / 6948x \\
Layers 8,12,16,20, independent & 128 & 1020 & 314.44 & 590.56 / 1.88x \\
Layers 8,12,16,20, independent & 256 & 1020 & 314.44 & 662.34 / 2.11x \\
Layers 8,16, independent & 256 & 4080 & 314.44 & 360.50 / 1.146x \\
Layers 8,16, joint logit-calibrated & 256 & 4080 & 314.44 & 308.76 / 0.982x \\
\bottomrule
\end{tabular}
\end{table}

\begin{table}[t]
\centering
\caption{Held-out and stress-test Qwen2.5-1.5B perplexity gates for patched
\mwg{} FFN layers on ASI3. Lower PPL ratio is better. The diverse64 suite is a
small mixed-domain stress set; heldout12 is a split held out from diverse
activation capture; Belle256 is a larger instruction-style validation slice.
The uncalibrated Belle256 rows use the first 50{,}838-token slice, while the
calibrated and persistent-baseline rows use the 51{,}267-token held-out slice
reserved for post-calibration evaluation.}
\label{tab:pplheldout}
\begin{tabular}{p{0.40\linewidth}rrrr}
\toprule
Patch policy & Rank & Tokens & Dense PPL & Patched PPL / ratio \\
\midrule
Layers 8,16, joint logit-calibrated, diverse64 & 256 & 5469 & 62.69 & 85.09 / 1.357x \\
Layer 16 only, diverse64 & 256 & 5469 & 62.69 & 69.24 / 1.105x \\
Layer 16 only, diverse64 & 384 & 5469 & 62.69 & 69.68 / 1.112x \\
Layer 16 only, heldout12 & 384 & 2560 & 58.14 & 61.71 / 1.061x \\
Layer 16 only, Belle256 & 256 & 50838 & 8.37 & 11.77 / 1.406x \\
Layer 16 only, Belle256 & 384 & 50838 & 8.37 & 10.77 / 1.287x \\
Layer 16 persistent low-rank, Belle256 & 384 & 51267 & 8.44 & 11.59 / 1.373x \\
Layer 16 only, Belle256, logit-MSE calibrated & 384 & 51267 & 8.44 & 10.03 / 1.189x \\
Layer 16 only, Belle256, LM-CE calibrated & 384 & 51267 & 8.44 & 9.96 / 1.180x \\
Layer 16 only, diverse64, LM-CE calibrated & 384 & 2736 & 62.37 & 67.80 / 1.087x \\
Layer 16 only, heldout12, LM-CE calibrated & 384 & 1280 & 58.12 & 60.48 / 1.040x \\
\bottomrule
\end{tabular}
\end{table}

\begin{table}[t]
\centering
\caption{Manifest-driven broad held-out validation for the ASI3 layer-16
rank-384 LM-CE checkpoint and same-rank persistent low-rank baselines. These
suites use explicit held-out files and reject empty or undersized text inputs.
Lower PPL ratio is better. The baseline was launched separately and has
different evaluated token counts, so it is a broad negative control rather than
a token-identical paired comparison.}
\label{tab:broadvalid}
\begin{tabular}{llrrrr}
\toprule
Policy & Suite & Tokens & Dense PPL & Patched PPL & Ratio \\
\midrule
MWG LM-CE & GSM8K test & 44{,}708 & 3.4941 & 4.7481 & 1.3589x \\
MWG LM-CE & MBPP test & 13{,}301 & 3.9877 & 4.5276 & 1.1354x \\
MWG LM-CE & Alpaca-tail & 33{,}343 & 3.8818 & 4.4249 & 1.1399x \\
\midrule
MWG aggregate & -- & 91{,}352 & -- & -- & 1.2464x \\
\midrule
Persistent LR & GSM8K test & 31{,}329 & 4.7059 & 6.0889 & 1.2939x \\
Persistent LR & MBPP test & 12{,}850 & 4.1360 & 5.3876 & 1.3026x \\
Persistent LR & Alpaca-tail & 22{,}694 & 4.2148 & 5.2535 & 1.2464x \\
\midrule
Persistent aggregate & -- & 66{,}873 & -- & -- & 1.2795x \\
\midrule
Persist+LMCE & GSM8K test & 44{,}708 & 3.4941 & 5.0680 & 1.4504x \\
Persist+LMCE & MBPP test & 13{,}301 & 3.9877 & 4.1225 & 1.0338x \\
Persist+LMCE & Alpaca-tail & 33{,}343 & 3.8818 & 4.7813 & 1.2317x \\
\midrule
Persist+LMCE agg. & -- & 91{,}352 & -- & -- & 1.3099x \\
\bottomrule
\end{tabular}
\end{table}

\begin{table}[t]
\centering
\caption{Token-level dense-fallback router validation on disjoint
router-train/router-eval text. The router is trained on token CE deltas and
then evaluated with actual mixed dense/patch forward passes, not offline oracle
loss mixing. Lower PPL ratio is better.}
\label{tab:tokenroutervalid}
\begin{tabular}{rrr}
\toprule
Target patch frac. & Actual token patch frac. & PPL ratio \\
\midrule
0.01 & 0.0038 & 1.0132x \\
0.03 & 0.0124 & 1.0249x \\
0.05 & 0.0193 & 1.0360x \\
0.10 & 0.0363 & 1.0508x \\
0.25 & 0.0947 & 1.1665x \\
0.50 & 0.2485 & 1.6013x \\
\bottomrule
\end{tabular}
\end{table}

\begin{table}[t]
\centering
\caption{Three-fold token-level router cross-validation over the pooled router
corpus. The 512 unique router lines are repartitioned into explicit 256/256
train/eval folds, then evaluated with actual mixed dense/patch forward passes.
This tests split sensitivity within the router distribution; it is not a
replacement for independent broad benchmarks. Lower PPL ratio is better.}
\label{tab:tokencvalid}
\begin{tabular}{rrrrr}
\toprule
Target patch frac. & Mean actual patch frac. & Mean ratio & Min ratio & Max ratio \\
\midrule
0.01 & 0.0009 & 0.9989x & 0.9955x & 1.0009x \\
0.03 & 0.0038 & 0.9940x & 0.9870x & 1.0009x \\
0.05 & 0.0075 & 0.9896x & 0.9796x & 1.0007x \\
0.10 & 0.0193 & 0.9787x & 0.9644x & 0.9974x \\
0.25 & 0.0695 & 0.9462x & 0.9282x & 0.9618x \\
0.50 & 0.1766 & 0.9218x & 0.9137x & 0.9373x \\
\bottomrule
\end{tabular}
\end{table}

\begin{table}[t]
\centering
\caption{Leakage-clean broad token-router validation. For each broad suite, the
router train pool removes exact lines overlapping that suite's eval file before
training. Values aggregate GSM8K, MBPP, and Alpaca-tail using token-weighted
losses. Lower PPL ratio is better.}
\label{tab:tokenbroadvalid}
\begin{tabular}{rrrrr}
\toprule
Target patch frac. & Actual patch frac. & Token-weighted ratio & Min suite & Max suite \\
\midrule
0.01 & 0.0005 & 1.0010x & 1.0001x & 1.0050x \\
0.03 & 0.0019 & 1.0008x & 1.0005x & 1.0025x \\
0.05 & 0.0055 & 0.9989x & 0.9970x & 1.0004x \\
0.10 & 0.0208 & 0.9916x & 0.9648x & 0.9996x \\
0.25 & 0.0849 & 0.9868x & 0.9197x & 1.0038x \\
0.50 & 0.1975 & 1.0140x & 0.8587x & 1.0498x \\
\bottomrule
\end{tabular}
\end{table}

\begin{table}[t]
\centering
\caption{Three-seed leakage-clean broad token-router sweep. Each seed uses
75\% of the overlap-clean router train candidates for each broad suite, then
evaluates actual mixed dense/patch forwards on GSM8K, MBPP, and Alpaca-tail.
Values report the mean across seeds; min and max are seed-level token-weighted
broad ratios. Lower PPL ratio is better.}
\label{tab:tokenbroadseeds}
\begin{tabular}{rrrrr}
\toprule
Target patch frac. & Mean actual patch frac. & Mean ratio & Min seed & Max seed \\
\midrule
0.01 & 0.0007 & 1.0011x & 1.0008x & 1.0013x \\
0.03 & 0.0021 & 0.9998x & 0.9980x & 1.0013x \\
0.05 & 0.0049 & 0.9977x & 0.9941x & 0.9997x \\
0.10 & 0.0182 & 0.9920x & 0.9890x & 0.9952x \\
0.25 & 0.0852 & 0.9864x & 0.9820x & 0.9929x \\
0.50 & 0.2168 & 1.0164x & 0.9943x & 1.0312x \\
\bottomrule
\end{tabular}
\end{table}

\begin{table}[t]
\centering
\caption{Leave-suite-out broad auxiliary token-router sweep. For each target
broad suite, router training uses the explicit router texts plus the other two
broad suites, with exact target-suite eval lines removed. Values report the
mean across three split seeds; min and max are seed-level token-weighted broad
ratios. Lower PPL ratio is better.}
\label{tab:tokenleavesuite}
\begin{tabular}{rrrrr}
\toprule
Target patch frac. & Mean actual patch frac. & Mean ratio & Min seed & Max seed \\
\midrule
0.01 & 0.0043 & 1.0007x & 1.0004x & 1.0009x \\
0.03 & 0.0166 & 0.9989x & 0.9980x & 1.0004x \\
0.05 & 0.0284 & 0.9975x & 0.9964x & 0.9996x \\
0.10 & 0.0565 & 0.9926x & 0.9914x & 0.9943x \\
0.25 & 0.1380 & 0.9943x & 0.9889x & 1.0004x \\
0.50 & 0.2806 & 1.0274x & 1.0076x & 1.0488x \\
\bottomrule
\end{tabular}
\end{table}

\begin{table}[t]
\centering
\caption{Deployable dense-fallback router validation on disjoint router
train/eval text. Values are averages over three launched ASI3 runs, which are
effectively identical because the current ridge router is deterministic for
these inputs. Lower PPL ratio is better.}
\label{tab:routervalid}
\begin{tabular}{lrrr}
\toprule
Policy & Target patch frac. & Actual token patch frac. & PPL ratio \\
\midrule
Routed frontier & 0.10 & 0.0169 & 1.0129x \\
Routed frontier & 0.25 & 0.0345 & 1.0273x \\
Routed frontier & 0.50 & 0.1453 & 1.1411x \\
Always patched & 1.00 & 1.0000 & 2.0625x \\
Oracle frontier & 0.10 & 0.1028 & 1.0139x \\
Oracle frontier & 0.25 & 0.2517 & 1.0624x \\
Oracle frontier & 0.50 & 0.5027 & 1.2000x \\
\bottomrule
\end{tabular}
\end{table}

\begin{table}[t]
\centering
\caption{ASI3 runtime profiler probe for the current unfused prototype. The
environment reported 8 visible Ascend 910B2 NPUs and available
\texttt{npu-smi}/\texttt{msprof}, but this PyTorch-profiler capture exposed
CPU activities only. It is therefore operator-level evidence, not full
off-chip hardware-counter validation. Lower time is better.}
\label{tab:runtimeprobe}
\begin{tabular}{lrrrr}
\toprule
Case & Descriptor MiB & Mean ms & Matmul-like ops & Fused-like op \\
\midrule
Dense & 64.50 & 0.3391 & 27 & no \\
\mwg{} $r=128$ & 5.53 & 0.7638 & 54 & no \\
\mwg{} $r=256$ & 11.06 & 0.7559 & 54 & no \\
\bottomrule
\end{tabular}
\end{table}

\paragraph{Interpretation.}
The most stable eight-rank point is $r=128$: it reduces descriptor traffic from
336.0 MiB to 13.5 MiB, reduces trainable FFN state from 336.0 MiB to
18.254 MiB after generator overhead, improves mean training latency from
$6.4874\pm0.8606$ ms to $4.2047\pm0.0299$ ms, and increases measured training
throughput from 1.279M to 1.948M tokens/s across eight ranks. The fastest
eight-rank point in this suite is $r=256$, reaching $1.572\pm0.165$x training
speedup and 1.992M tokens/s, while still reducing descriptor traffic by 12.4x
and trainable FFN state by 10.3x. The lower-rank $r=64$ point gives the
strongest compression, with 49.8x descriptor-traffic reduction and 30.2x
trainable-state reduction. Across 2, 4, and 8 ranks, the all-reduce proxy drops
from dense latencies of 19.66, 11.49, and 6.91 ms to low-rank latencies near
1--2.5 ms, confirming that the trainable-state reduction translates into a
distributed communication benefit. The forward-only path remains sensitive to
unfused factor-matmul overhead; fused descriptor generation, on-chip factor
consumption, and no descriptor write-back remain the systems requirements that
separate \mwg{} from an ordinary low-rank module.

Table \ref{tab:quality} is deliberately more conservative. On a real
Qwen2.5-1.5B FFN layer, static SVD at $r=64$ and $r=128$ is a weak
approximation despite strong byte savings. Short distillation improves relative
MSE and cosine, but the current rank-scale \mwg{} generator is almost tied with
the persistent low-rank baseline. This result prevents an overclaim: the
present prototype establishes the communication and latency opportunity, while
future work must use stronger conditional generators, longer distillation, or
end-to-end training to convert that systems opportunity into language quality.

Tables \ref{tab:ppldiag}, \ref{tab:pplheldout}, \ref{tab:broadvalid},
\ref{tab:tokenroutervalid}, \ref{tab:tokencvalid},
\ref{tab:tokenbroadvalid}, \ref{tab:tokenbroadseeds},
\ref{tab:tokenleavesuite}, and
\ref{tab:routervalid} show the resulting pivot, while Table
\ref{tab:runtimeprobe} records the current runtime instrumentation boundary.
Early-layer replacement is not
credible: even rank-384 experiments left layer 0 with unacceptable perplexity
damage, and replacing layers 0, 4, 8, and 12 at rank 128 collapses perplexity.
The usable regime is instead quality-gated and patch-aware, but the latest
ASI3 stress test also shows how easy it is to overfit a small gate. Joint
logit calibration of layers 8 and 16 reaches a 0.982x PPL ratio on the
4080-token repeated-text gate, but degrades to 1.357x on the diverse64 stress
set. The more credible current policy is therefore stricter: keep early layers,
layer 4, layer 8, layer 24, and multi-layer compositions dense unless they pass
held-out gates, and patch only the safest eligible layer. In these experiments,
layer 16 is the best candidate: rank 256 gives a 1.105x PPL ratio on diverse64,
while a rank-384 retraining on diverse activation data gives 1.061x on the
held-out split. However, the same layer-16 policy is not yet robust: on a
larger Belle-style instruction slice, rank 384 has a 1.287x PPL ratio before
logit calibration. End-to-end logit-MSE calibration improves that
external-corpus gate monotonically in the current sweep, reaching 1.205x after
400 calibration steps and 1.189x after longer gentle continuation. Direct
language-model cross-entropy (LM-CE) calibration is the stronger current
variant: it reduces the Belle held-out ratio to 1.180x and improves the small
off-distribution checks to 1.087x on diverse64 and 1.040x on heldout12. A
same-rank persistent low-rank baseline trained on the same layer-16 activation
cache reaches 1.373x on the Belle held-out gate, so the calibrated MWG variant
is better than this simple durable low-rank control in the present setting.
However, the remaining degradation is still material, the diverse checks are
small, and the result is from one checkpoint and one split. Table
\ref{tab:broadvalid} is the decisive current counterweight: when the same
layer-16 checkpoint is evaluated on explicit GSM8K, MBPP, and Alpaca-tail
held-out files, all three suites degrade, with a 1.246x token-weighted aggregate
PPL ratio. A same-rank persistent low-rank baseline trained from explicit
router-train activations also fails the broad manifest, with a 1.279x
token-weighted ratio in its separate launch; adding direct LM-CE calibration to
that persistent baseline worsens the broad aggregate to 1.310x over the
91{,}352-token manifest. Thus the calibrated \mwg{} variant
is not merely losing to a trivial durable low-rank baseline, but neither method
is close to broad quality preservation. Tables \ref{tab:tokenroutervalid} and
\ref{tab:routervalid} show that dense fallback can bound quality loss only by
patching very few tokens. The stricter token-level router is trained on token
CE deltas and then evaluated with actual mixed dense/patch forward passes. It
exposes a large split shift: on router-train the patch is favorable
(``patched'' PPL ratio 0.437x over 30{,}053 tokens), but on disjoint
router-eval it is unfavorable (2.062x over 35{,}118 tokens). Mixed forwarding
keeps ratios near 1.01--1.05x only while patching about 0.4--3.6\% of tokens,
then degrades to 1.166x at 9.5\% patching and 1.601x at 24.8\% patching.
However, Table \ref{tab:tokencvalid} shows that this result is split-sensitive
within the router distribution: when the two router corpora are pooled and
repartitioned into three explicit train/eval folds, mean mixed-forward ratios
fall below dense, reaching 0.946x at 6.9\% actual patching and 0.922x at
17.7\% patching. This is an encouraging signal for risk-aware routing, but it
does not by itself overturn the negative broad manifest because the folds are
resamples of the router corpus rather than independent GSM8K, MBPP, or
Alpaca-style benchmarks. Table \ref{tab:tokenbroadvalid} is the stricter check:
for each broad suite, exact eval-line overlaps are removed from the router
train pool before fitting the token router. Always-patched remains negative on
all suites (1.324x on GSM8K, 1.113x on MBPP, and 1.110x on Alpaca-tail), so the
patch itself is not broadly quality preserving. Selective routing is better:
the broad token-weighted frontier reaches 0.992x at 2.1\% actual patching and
0.987x at 8.5\% actual patching, before degrading to 1.014x at 19.8\% patching.
Table \ref{tab:tokenbroadseeds} repeats this leakage-clean broad protocol over
three explicit split seeds using 75\% of each overlap-clean train pool. The
bounded selective result repeats: target 0.05, 0.10, and 0.25 are below dense on
every seed, with mean token-weighted ratios of 0.998x, 0.992x, and 0.986x,
respectively. The high-patching target 0.50 row is unstable, with seed-level
ratios spanning 0.994x to 1.031x. This supports a bounded selective-routing
claim rather than an always-patched deployment claim. Table
\ref{tab:tokenleavesuite} is a stricter distribution-robust check: for each
target broad suite, the router also trains on the two non-target broad suites
but removes exact target-suite eval lines. The bounded signal survives at target
0.05 and 0.10 on every seed, with mean ratios of 0.997x and 0.993x. Target 0.25
is below dense on average but has one near-neutral/slightly-worse seed, and
target 0.50 degrades on every seed. The example-level routed frontier similarly
keeps ratios near
1.01--1.03x at 1.7--3.5\% actual token patching, reaches 1.141x by 14.5\%
patching, and always-patched evaluation degrades to 2.063x on the router-eval
split. The oracle frontier is better at the same nominal target fractions but
is not deployable because it uses evaluation losses. At the present stage this
is a reproducible feasibility signal plus a bounded selective-router result,
not yet a top-journal-quality deployment claim; stronger training, stronger
baselines, broader independent benchmark families, and hardware-counter or
fused-kernel evidence remain required. The runtime probe in Table
\ref{tab:runtimeprobe} sharpens that last limitation: even on an ASI3 host with
visible Ascend 910B2 NPUs and available \texttt{npu-smi}/\texttt{msprof}, the
current unfused associative prototype appears as ordinary matmul-like
operators, doubles the matmul-like op count in this microprobe, and is slower
than dense. It is useful negative fused-kernel evidence and tool-availability
evidence, but it is not yet an off-chip traffic counter validation.

\section{Discussion}

\mwg{} is most useful when FFN weight traffic dominates and when the target
hardware has enough arithmetic headroom to spend extra low-rank multiplies or
generator work. It is less attractive when the model is already compute-bound,
when rank must approach full dimension, or when the generator is placed so far
from the FFN kernel that descriptor traffic reappears.

The architecture also changes the meaning of compression. The goal is not only
to reduce checkpoint size, but to change where parameters live during execution.
That distinction is why the method belongs partly to neural architecture design
and partly to model-runtime infrastructure.

\paragraph{Relation to KV-cache compression and MLA.}
The KV-cache memory used during autoregressive decoding is a different object
from the FFN weight traffic studied here. A standard attention layer stores
past keys and values with leading size $\Theta(2L n_h d_h)$ elements for
sequence length $L$, head count $n_h$, and head dimension $d_h$; in practical
long-context serving this footprint is usually an off-chip memory allocation,
although recently accessed tiles may pass through on-chip caches during an
attention kernel. MLA reduces this attention-side storage by caching a latent
key-value vector and reconstructing or absorbing key/value projections at
attention time \cite{deepseek2024v2}. \mwg{} is orthogonal: it targets FFN
projection weights, does not by itself compress the attention KV cache, and
uses low-rank generation to change the lifecycle of weight operands rather
than the representation of past-token activations.

\paragraph{Current limitations.}
Three gaps are intentionally left visible. First, the end-to-end quality
evidence is still negative outside the narrowest gates: it covers a small
Qwen2.5-1.5B text gate, a stricter single-layer eligibility policy, one larger
instruction-style stress test that still shows material degradation, and a
new broad held-out validation where the same layer-16 checkpoint degrades
GSM8K, MBPP, and Alpaca-tail. The added manifest harness prevents default-text
self-evaluation, and the deployable-router harness separates train-time risk
labels from held-out routing decisions, but the current router preserves dense
quality only at very small patch-token fractions. Second, descriptor lifecycle
claims should be confirmed with low-level hardware performance counters for
off-chip read and write transactions, not only with software-side theoretical
byte estimates. Third, the current generator and associative path are
deliberately small and benchmarkable; stronger generators and fused kernels may
improve both quality and forward latency, but must be measured against their own
placement and traffic costs.

\section{Conclusion}

This paper reframes Transformer FFN weights as generated ephemeral operands.
By replacing persistent dense matrices with conditional rank-$r$ descriptors and
an associative execution path, \mwg{} reduces both memory traffic and
distributed gradient communication in an eight-rank 8B-scale FFN systems
benchmark. A checkpoint-backed Qwen2.5-1.5B probe adds the missing negative
control: static and short-distilled low-rank FFNs do not recover dense quality
when applied indiscriminately, especially in early layers or brittle multi-layer
compositions. The stronger result is a stricter quality-gated policy: layer 16
alone is the best current candidate on small diverse gates, whereas a two-layer
calibrated patch overfits the original small gate and the layer-16 patch still
degrades on a larger instruction-style corpus, although LM-CE calibration
reduces that degradation and improves the small diverse checks. The accompanying
benchmark is therefore best read as a reproducible systems case for ephemeral
operand lifecycles plus an initial algorithmic recipe and negative evidence for
quality-preserving layer eligibility. Broad dataset validation now shows that
the present checkpoint is insufficient; stronger training, stronger routing,
and fused-kernel counter evidence are the next experimental steps.

\bibliographystyle{plain}
\bibliography{refs}

\end{document}
